\documentclass[11pt,a4paper]{article}
\usepackage{get-uebung}

% Knotennummer im Kreis: \knoten{(x,y)}{Nr}
\newcommand{\knoten}[2]{\node[draw,circle,inner sep=1pt,font=\scriptsize,hsublau] at #1 {#2};}

\blatt{11}{Knotenpotentialverfahren}{Kap.~2.14}

\begin{document}

\blattkopf{
  \item das Kochrezept des Knotenpotentialverfahrens Schritt für Schritt anwenden,
  \item reale Spannungsquellen in reale Stromquellen umwandeln und Widerstände als Leitwerte schreiben,
  \item Leitwertmatrix und rechte Seite direkt aus der Schaltung ablesen,
  \item ein $2\times2$-Gleichungssystem von Hand lösen und daraus alle Zweigströme bestimmen,
  \item ideale Spannungsquellen durch Verschieben durch einen Knoten behandeln.
}

% -------------------------------------------------------------------
\begin{aufgabe}{Gleichungssystem direkt aufstellen}{15}
Die Knoten der Schaltung sind bereits nummeriert, Knoten \textcircled{\scriptsize 0} ist der Bezugsknoten.
\begin{center}
\begin{circuitikz}
  \draw (0,0) to[I, i=$I_1$] (0,3) -- (2,3);
  \draw (2,3) to[R, l_=$G_1$, *-*] (2,0);
  \draw (2,3) to[R, l=$G_2$] (5,3);
  \draw (5,3) to[R, l=$G_3$, *-*] (5,0);
  \draw (5,3) to[R, l=$G_4$, *-] (8,3);
  \draw (8,3) to[R, l=$G_5$, *-*] (8,0);
  \draw (2,3) -- (2,5.4) to[R, l=$G_6$] (8,5.4) -- (8,3);
  \draw (8,4.2) to[I, i^=$I_2$, *-] (5,4.2) -- (5,3);
  \draw (0,0) -- (8,0);
  \knoten{(1.6,3.4)}{1}
  \knoten{(4.6,3.4)}{2}
  \knoten{(8.45,3.4)}{3}
  \knoten{(4.6,-0.4)}{0}
\end{circuitikz}
\end{center}
\begin{teilaufgaben}
  \item Stellen Sie mit dem Kochrezept das Gleichungssystem $\mathbf{G}\cdot\vec{U} = \vec{I}$ für die Knotenpotentiale $U_{10}$, $U_{20}$, $U_{30}$ auf. (Nicht lösen!)
  \item Die Leitwertmatrix ist symmetrisch. Woran liegt das?
  \item Leiten Sie die erste Zeile des Gleichungssystems zur Kontrolle mit dem Knotensatz an Knoten \textcircled{\scriptsize 1} her.
  \item Wie ändert sich das Gleichungssystem, wenn die Stromquelle $I_2$ umgedreht wird? Und wie, wenn $G_6$ entfernt wird?
\end{teilaufgaben}
\end{aufgabe}

\begin{loesung}
\schritt{Vorbereitung}
Knoten sind nummeriert, es gibt keine Spannungsquellen, alle Elemente sind schon Leitwerte -- wir können sofort aufstellen.
\schritt{a) Leitwertmatrix}
\leitfrage{Welche Leitwerte sind direkt mit Knoten 1 verbunden? Welche verbinden Knoten 1 direkt mit Knoten 2 bzw.\ 3?}
\begin{itemize}[nosep]
  \item Hauptdiagonale $G_{ii}$: Summe \emph{aller} Leitwerte am Knoten $i$:\\
        $G_{11} = G_1+G_2+G_6$, \quad $G_{22} = G_2+G_3+G_4$, \quad $G_{33} = G_4+G_5+G_6$
  \item Nebenelemente $G_{ij}$: \emph{negative} Summe der Leitwerte zwischen $i$ und $j$:\\
        $G_{12} = G_{21} = -G_2$, \quad $G_{13} = G_{31} = -G_6$, \quad $G_{23} = G_{32} = -G_4$
\end{itemize}
\schritt{a) Rechte Seite}
Quellenströme \emph{in} den Knoten positiv, \emph{heraus} negativ: $I_1$ fließt in Knoten 1; $I_2$ fließt aus Knoten 3 heraus und in Knoten 2 hinein.
\[
  \begin{pmatrix}
    G_1+G_2+G_6 & -G_2 & -G_6\\
    -G_2 & G_2+G_3+G_4 & -G_4\\
    -G_6 & -G_4 & G_4+G_5+G_6
  \end{pmatrix}
  \cdot
  \begin{pmatrix} U_{10}\\ U_{20}\\ U_{30}\end{pmatrix}
  = \ergebnis{\begin{pmatrix} I_1\\ I_2\\ -I_2\end{pmatrix}}
\]
\schritt{b) Symmetrie}
Ein Leitwert zwischen den Knoten $i$ und $j$ verbindet $j$ genauso mit $i$. Deshalb ist $G_{ij} = G_{ji}$. Ideale Stromquellen tauchen nie in der Matrix auf, nur auf der rechten Seite.
\schritt{c) Kontrolle mit dem Knotensatz}
Ströme, die Knoten 1 über die Leitwerte verlassen, sind gleich dem zufließenden Quellenstrom:
\[
  G_1U_{10} + G_2(U_{10}-U_{20}) + G_6(U_{10}-U_{30}) = I_1
  \;\Leftrightarrow\; (G_1+G_2+G_6)U_{10} - G_2U_{20} - G_6U_{30} = I_1 \;\checkmark
\]
\schritt{d) Änderungen}
$I_2$ umgedreht: Nur die rechte Seite ändert sich zu $(I_1,\,-I_2,\,+I_2)^T$, die Matrix bleibt.\\
$G_6$ entfernt: $G_6$ verschwindet aus $G_{11}$ und $G_{33}$, und $G_{13} = G_{31} = 0$.
\typfehler{Nebenelemente mit positivem Vorzeichen eintragen. Die Leitwerte der Stromquellenzweige in die Matrix aufnehmen. Beim Quellenstrom zwischen zwei Knoten nur einen der beiden Knoten berücksichtigen -- $I_2$ taucht immer in zwei Zeilen mit entgegengesetztem Vorzeichen auf.}
\end{loesung}

% -------------------------------------------------------------------
\begin{aufgabe}{Knotenpotentialverfahren mit Zahlen}{22}
\begin{center}
\begin{circuitikz}
  \draw (0,0) to[V, v<=$U_1$] (0,3) to[R, l=$R_1$] (3,3);
  \draw (3,3) to[R, l_=$R_2$, *-*] (3,0);
  \draw (3,3) to[R, l=$R_3$] (6,3);
  \draw (6,3) to[R, l_=$R_4$, *-*] (6,0);
  \draw (6,3) -- (8,3); \draw (8,0) to[I, i=$I_q$] (8,3); \draw (8,0) -- (0,0);
\end{circuitikz}
\end{center}
\gegeben{$U_1 = \SI{10}{V}$, $I_q = \SI{1.5}{A}$, $R_1 = \SI{2}{\ohm}$, $R_2 = \SI{4}{\ohm}$, $R_3 = \SI{4}{\ohm}$, $R_4 = \SI{2}{\ohm}$}
\begin{teilaufgaben}
  \item Nummerieren Sie die Knoten (Bezugsknoten unten) und wandeln Sie die reale Spannungsquelle ($U_1$, $R_1$) in eine reale Stromquelle um. Zeichnen Sie die umgewandelte Schaltung mit Leitwerten.
  \item Stellen Sie das Gleichungssystem für die Knotenpotentiale auf.
  \item Lösen Sie das Gleichungssystem von Hand (z.\,B.\ mit der Cramerschen Regel).
  \item Berechnen Sie die Ströme durch alle vier Widerstände der \emph{Originalschaltung} und prüfen Sie den Knotensatz an beiden Knoten.
\end{teilaufgaben}
\begin{tipp}Für $\begin{pmatrix}a&b\\c&d\end{pmatrix}\vec{x} = \begin{pmatrix}e\\f\end{pmatrix}$ gilt $x_1 = \dfrac{ed-bf}{ad-bc}$, $x_2 = \dfrac{af-ec}{ad-bc}$.\end{tipp}
\end{aufgabe}

\begin{loesung}
\schritt{a) Vorbereitung}
Knoten \textcircled{\scriptsize 1}: oberer Anschluss von $R_2$; Knoten \textcircled{\scriptsize 2}: oberer Anschluss von $R_4$; Knoten \textcircled{\scriptsize 0}: unten.\\
Quellenumwandlung (Innenwiderstand bleibt, Pfeil der Stromquelle zeigt zum Knoten 1):
\[
  I_{q1} = \frac{U_1}{R_1} = \frac{\SI{10}{V}}{\SI{2}{\ohm}} = \SI{5}{A}, \qquad G_1 = \frac{1}{R_1} = \SI{0.5}{S}
\]
Leitwerte: $G_2 = \SI{0.25}{S}$, $G_3 = \SI{0.25}{S}$, $G_4 = \SI{0.5}{S}$.
\begin{center}
\begin{circuitikz}
  \draw (0,0) to[I, i=$I_{q1}$] (0,3) -- (3,3);
  \draw (1.5,3) to[R, l_=$G_1$, *-*] (1.5,0);
  \draw (3,3) to[R, l=$G_2$, *-*] (3,0);
  \draw (3,3) to[R, l=$G_3$] (6,3);
  \draw (6,3) to[R, l_=$G_4$, *-*] (6,0);
  \draw (6,3) -- (8,3); \draw (8,0) to[I, i=$I_q$] (8,3); \draw (8,0) -- (0,0);
  \knoten{(3,3.45)}{1}
  \knoten{(6,3.45)}{2}
  \knoten{(4.5,-0.4)}{0}
\end{circuitikz}
\end{center}
\schritt{b) Gleichungssystem}
\[
  \begin{pmatrix} G_1+G_2+G_3 & -G_3\\ -G_3 & G_3+G_4\end{pmatrix}
  \begin{pmatrix} U_{10}\\U_{20}\end{pmatrix}
  = \begin{pmatrix} I_{q1}\\ I_q\end{pmatrix}
  \quad\Rightarrow\quad
  \begin{pmatrix} 1 & -0{,}25\\ -0{,}25 & 0{,}75\end{pmatrix}\si{S}
  \begin{pmatrix} U_{10}\\U_{20}\end{pmatrix}
  = \begin{pmatrix} 5\\ 1{,}5\end{pmatrix}\si{A}
\]
\schritt{c) Lösen}
Determinante: $D = 1\cdot 0{,}75 - (-0{,}25)^2 = 0{,}6875\ \si{S^2}$.
\[
  U_{10} = \frac{5\cdot 0{,}75 - (-0{,}25)\cdot 1{,}5}{0{,}6875}\,\si{V} = \frac{4{,}125}{0{,}6875}\,\si{V} = \ergebnis{\SI{6}{V}}
\]
\[
  U_{20} = \frac{1\cdot 1{,}5 - 5\cdot(-0{,}25)}{0{,}6875}\,\si{V} = \frac{2{,}75}{0{,}6875}\,\si{V} = \ergebnis{\SI{4}{V}}
\]
\schritt{d) Zweigströme in der Originalschaltung}
\leitfrage{Darf man den Strom durch $R_1$ als $G_1\cdot U_{10}$ berechnen?}
Nein! Durch $G_1$ in der umgewandelten Schaltung fließt ein anderer Strom als durch $R_1$ in der Originalschaltung. In der Originalschaltung liegt an $R_1$ die Spannung $U_1 - U_{10}$:
\begin{align*}
  I_{R1} &= \frac{U_1-U_{10}}{R_1} = \frac{\SI{4}{V}}{\SI{2}{\ohm}} = \ergebnis{\SI{2}{A}} \text{ (zu Knoten 1)}, &
  I_{R2} &= \frac{U_{10}}{R_2} = \ergebnis{\SI{1.5}{A}},\\
  I_{R3} &= \frac{U_{10}-U_{20}}{R_3} = \ergebnis{\SI{0.5}{A}} \text{ (von 1 nach 2)}, &
  I_{R4} &= \frac{U_{20}}{R_4} = \ergebnis{\SI{2}{A}}.
\end{align*}
Knoten 1: $\SI{2}{A} = \SI{1.5}{A} + \SI{0.5}{A}$ \checkmark \qquad Knoten 2: $\SI{0.5}{A} + \SI{1.5}{A} = \SI{2}{A}$ \checkmark\\
(Zum Vergleich: In der umgewandelten Schaltung fließen $G_1U_{10} = \SI{3}{A}$ durch $G_1$, und $\SI{5}{A}-\SI{3}{A} = \SI{2}{A}$ in den Knoten -- an den Klemmen verhalten sich beide Quellen gleich, intern nicht.)
\typfehler{Die umgewandelte Stromquelle mit falscher Richtung einzeichnen (ihr Pfeil zeigt dorthin, wo der Pluspol der Spannungsquelle war). Den Strom durch $R_1$ aus der umgewandelten Schaltung ablesen.}
\end{loesung}

% -------------------------------------------------------------------
\begin{aufgabe}{Ideale Spannungsquelle}{20}
\begin{center}
\begin{circuitikz}
  \draw (0,0) to[V, v<=$U$] (0,4) -- (6,4);
  \draw (2,4) to[R, l_=$R_1$, *-*] (2,2) to[R, l_=$R_2$, -*] (2,0);
  \draw (6,4) to[R, l=$R_3$] (6,2) to[R, l=$R_4$, *-] (6,0) -- (0,0);
  \draw (2,2) to[R, l=$R_5$] (6,2);
  \knoten{(1.5,2.4)}{1}
  \knoten{(6.5,2.4)}{2}
  \knoten{(4,4.4)}{3}
  \knoten{(4,-0.4)}{0}
\end{circuitikz}
\end{center}
\gegeben{$U = \SI{12}{V}$, $R_1 = R_4 = \SI{2}{\ohm}$, $R_2 = R_3 = R_5 = \SI{4}{\ohm}$}
\begin{teilaufgaben}
  \item Warum lässt sich das Kochrezept hier nicht direkt anwenden?
  \item Verschieben Sie die ideale Spannungsquelle durch den Knoten \textcircled{\scriptsize 3} in die Zweige von $R_1$ und $R_3$. Zeichnen Sie die neue Schaltung. Was passiert mit Knoten \textcircled{\scriptsize 3}?
  \item Stellen Sie das Gleichungssystem für $U_{10}$ und $U_{20}$ auf und lösen Sie es.
  \item Berechnen Sie den Strom, den die Spannungsquelle abgibt.
\end{teilaufgaben}
\end{aufgabe}

\begin{loesung}
\schritt{a) Problem}
Eine \emph{ideale} Spannungsquelle hat keinen Innenwiderstand. Man kann sie deshalb nicht in eine Stromquelle umwandeln ($I_k = U/0 \to \infty$).
\schritt{b) Quelle verschieben}
\leitfrage{Welche Spannung liegt zwischen Knoten 3 und den oberen Enden von $R_1$ und $R_3$?}
Wir schieben die Quelle durch Knoten 3 hindurch und setzen je eine Kopie in Reihe zu $R_1$ und zu $R_3$ (Skript Abb.~2.7). Die oberen Enden von $R_1$ und $R_3$ liegen weiterhin um $U$ über dem Bezugsknoten, an den Spannungen ändert sich nichts.
Knoten 3 liegt jetzt auf \SI{0}{V} und fällt mit dem Bezugsknoten zusammen:
\begin{center}
\begin{circuitikz}
  \draw (0,0) to[V, v<=$U$] (0,2) to[R, l=$R_1$] (2,2);
  \draw (2,2) to[R, l_=$R_2$, *-*] (2,0);
  \draw (2,2) to[R, l=$R_5$] (6,2);
  \draw (6,2) to[R, l=$R_4$, *-*] (6,0);
  \draw (8,0) to[V, v_<=$U$] (8,2) to[R, l_=$R_3$] (6,2);
  \draw (0,0) -- (8,0);
  \knoten{(2,2.45)}{1}
  \knoten{(6,2.45)}{2}
  \knoten{(4,-0.4)}{0}
\end{circuitikz}
\end{center}
Jetzt sind ($U$, $R_1$) und ($U$, $R_3$) reale Spannungsquellen und lassen sich umwandeln:
$\;U/R_1 = \SI{6}{A}$ parallel zu $G_1 = \SI{0.5}{S}$ in Knoten 1, $\;U/R_3 = \SI{3}{A}$ parallel zu $G_3 = \SI{0.25}{S}$ in Knoten 2.
\schritt{c) Gleichungssystem und Lösung}
Mit $G_2 = G_5 = \SI{0.25}{S}$, $G_4 = \SI{0.5}{S}$:
\[
  \begin{pmatrix} G_1+G_2+G_5 & -G_5\\ -G_5 & G_3+G_4+G_5\end{pmatrix}
  \begin{pmatrix} U_{10}\\U_{20}\end{pmatrix}
  = \begin{pmatrix} UG_1\\ UG_3\end{pmatrix}
  \;\Rightarrow\;
  \begin{pmatrix} 1 & -0{,}25\\ -0{,}25 & 1\end{pmatrix}\si{S}
  \begin{pmatrix} U_{10}\\U_{20}\end{pmatrix}
  = \begin{pmatrix} 6\\ 3\end{pmatrix}\si{A}
\]
$D = 1 - 0{,}0625 = 0{,}9375\ \si{S^2}$:
\[
  U_{10} = \frac{6\cdot 1 + 0{,}25\cdot 3}{0{,}9375}\,\si{V} = \ergebnis{\SI{7.2}{V}},\qquad
  U_{20} = \frac{1\cdot 3 + 6\cdot 0{,}25}{0{,}9375}\,\si{V} = \ergebnis{\SI{4.8}{V}}
\]
\schritt{d) Quellenstrom}
Zurück in der \emph{Originalschaltung}: Die Quelle speist die Zweige $R_1$ und $R_3$, Knoten 3 liegt auf $U = \SI{12}{V}$:
\[
  I_U = \frac{U-U_{10}}{R_1} + \frac{U-U_{20}}{R_3} = \frac{\SI{4.8}{V}}{\SI{2}{\ohm}} + \frac{\SI{7.2}{V}}{\SI{4}{\ohm}} = \SI{2.4}{A}+\SI{1.8}{A} = \ergebnis{\SI{4.2}{A}}
\]
Kontrolle Knoten 1: zu $\SI{2.4}{A}$, ab $U_{10}/R_2 + (U_{10}-U_{20})/R_5 = \SI{1.8}{A}+\SI{0.6}{A}$ \checkmark
\typfehler{Die Quelle nur in \emph{einen} der Zweige verschieben -- sie muss in \emph{jeden} Zweig, der am Knoten hängt. Vergessen, dass Knoten 3 nach dem Verschieben mit dem Bezugsknoten zusammenfällt, und ihn weiter als Unbekannte führen.}
\end{loesung}

\end{document}
